
\section{Pseudo-quadratic-like maps and pseudo-puzzle}\label{pseudo-puzzle}

In this section, we will summarize the needed background on pseudo-quadratic-like maps
and pseudo-puzzle. The details can be found in \cite{K,decorations}. 

\subsection{Pseudo-quadratic-like maps}

Suppose that $\BU'$, $\BU$ are disks,
$i:  \BU' \to \BU$ is a holomorphic immersion,
and $f: \BU' \to \BU$ is a degree $d$ holomorphic branched cover.
Suppose further that there exist full continua $K \Subset \BU$ and $K'\Subset \BU'$
such that $K' = i^{-1}(K) = f^{-1}(K)$.
Then we say that $F=(i,f):  \BU'\ra (\BU, \BU)$  is a  
$\psi$-{\it quadratic-like}
 ($\psi$-{\it ql}) map  with filled Julia set $K$.
We let
$$
   \psimod(F)= \psimod(f) = \mod(\BU\sm K). 
$$ 

\begin{lem}\label{ql extension}
Let $F=(i,f)\from\BU'\ra \BU$ be a $\psi$-ql map of degree $d$ 
 with filled Julia set $K$.
Then $i$ is an embedding in  a neighborhood of $K' \equiv f^{-1}(K)$, 
and the map  $g\equiv  f\circ i^{-1}\from U'\ra U$ near $K$ is quadratic-like.

\ssk Moreover, the domains $U$ and $U'$ can be selected in such a way that
$$ \mod(U\sm i(U'))\geq \mu (\psimod(F) >0.$$ 
\end{lem}

There is a natural $\psi$-ql map $\BU^n\ra \BU^{n-1}$, 
the ``restriction'' of $(i,f)$ to $\BU^n$. 
Somewhat loosely, we will use the same notation $F=(i,f)$ for this restriction.  

 Let us normalize  the $\psi$-quadratic-like  maps under consideration 
so that $\diam K'= \diam K=1$, both $K$ and $K'$ contain 0 and 1,  
$0$ is the critical point of $f$, and $i(0)=0$.
Let us endow the space of $\psi$-quadratic-like maps (considered up to independent rescalings in the domain and the range)
with the Carath\'eodory topology.
In this topology, a sequence of normalized maps $(i_n, f_n): \BU_n'\ra\BU_n$ converges to $(i,f): \BU'\ra \BU$ if
the pointed domains $(\BU_n', 0)$ and $(\BU_n, 0)$ converge to $\BU'$ and $\BU$ respectively,
and the maps $i_n$, $f_n$ converge respectively to $i$, $f$, 
uniformly on compact subsets of  $\BU'$.  
  

\begin{lem}\label{compactness}
  Let $\mu>0$. Then the space of $\psi$-PL maps $F$ with connected  Julia set
and  $\psimod(F)\geq \mu$ is compact.
\end{lem}  

To simplify notation, we will often refer to $f$ as a ``$\psi$-ql map'' 
keeping $i$ in mind implicitly. 


\subsection{Pseudo-puzzle}\label{psi-puzzle}

\subsubsection{Definitions}
   Let $(i,f) : \BU'\ra \BU$ be a $\psi$-ql map.
By Lemma~\ref{ql extension}, it admits a quadratic-like restriction $U'\ra U$ to a neighborhood of
its (filled) Julia set  $K=K_\BU$.
Here $U'$ is embedded into $U$, so we can identify $U'$ with $i(U')$
and $ f: U'\ra U $ with  $f \circ i^{-1} $. 

Assume that $K$ is connected and both fixed points of $f$ are repelling.
Then we can cut $U$ by external rays landing at the $\alpha$-fixed point and consider the corresponding Yoccoz puzzle.

Given a (geometric) puzzle piece $Y$ of bidepth $(m,l)$,
recall that $K_Y$ stands for $Y\cap K(f)$. % and  $S_Y = \cl (K(f) \sm K_Y)$. 
Let us consider the topological annulus $A=\BU^l  \sm K(f)$ and its universal covering $\hat A$. 
Let $Y_i$ be the components of $Y\sm K_Y$. There are finitely many of them,
and each $Y_i$ is simply connected. Hence they can be embedded into $\hat A$.
Select such an embedding $e_i: Y_i\ra \hat A_i$ where $\hat A_i$ stands for a copy of $\hat A$. 
Then glue the $A_i$ to $Y$ by means of $e_i$, i.e., let $\BY= Y \sqcup_{e_i} \hat A_i.$
This is  the {\it pseudo-piece} (``$\psi$-piece'') associated with $Y$.
Noter that the Julia piece $K_Y$ naturally embeds into $\BY$.

\begin{lem}\label{psi-maps}
\begin{itemize}
\item [(i)]  Consider two puzzle pieces $Y$ and $Z$ such that the map $f: Y\ra Z$ is a branched covering
of degree $k$ (where $k=1$ or $k=2$ depending on whether $Y$ is off-critical or not).
% Then  the inverse map $f^{-1}: Z\ra Y$ extends to a $k$-valued immersion
Then there exists an induced map $\f: \BY \ra \BZ$
which  is a branched covering of the same degree $k$.  


\item [(ii)] Given two puzzle pieces $Y\subset Z$, the inclusion $i: Y\ra Z$ extends to an immersion $\i: \BY\ra \BZ$.
\end{itemize}
\end{lem}

\comm{******************************************
\sss{Moduli}
 Given two puzzle pieces $Z\Subset Y$, we let 
$$
  \psimod (Y,Z) = \mod(\BY\sm K_Z).
$$

Lemma \ref{psi-maps} implies:

\begin{lem}\label{psi-mod}
\begin{itemize}
  \item [(i)]  Consider two pairs of puzzle pieces $(Y',Y)$ and $(Z',Z)$ 
such that the map $f: (Y',Y)\ra (Z',Z)$ is a branched covering
of degree $k$ (on both domains).
Then $$\psimod(Z',Z)=k\, \psimod(Y',Y).$$

\item [(ii)] Given a nest of three puzzle pieces $W\subset Z\subset  Y$, 
we have $$\psimod(Z,W)\leq \psimod(Y,W).$$
\end{itemize}
\end{lem}
*****************************************************}

\subsubsection{Boundary of puzzle pieces}\label{boundary}
The ideal boundary of a $\psi$-puzzle piece $\BY$
is tiled by (finitely many) arcs $\lambda_i\subset \di \hat A_i$
that cover the ideal boundary of $\BU^m$ (where $m=\depth Y$) 
and arcs $\xi_i, \eta_i \subset \di \hat A_i$ mapped onto the Julia set $J(f)$. 
The arc $\la_i$ meets each $\xi_i,\, \eta_i$ at a single boundary point
corresponding to a path $\de: [0,1)\mapsto A$ that wraps around $K(f)$ infinitely many times,
while $\eta_i$ meets $\xi_{i+1}$ at a vertex $v_i\in Y\cap K(f)$. 
We say  that the arcs $\la_i$ form the {\it outer boundary } (or ``$O$-boundary'')
 $\di_O\BY$ of the puzzle piece $\BY$,
while the arcs $\xi_i$ and $\eta_i$
form its {\it $J$-boundary} $\di_J \BY$. 
Given a vertex $v=v_i$ of a puzzle piece $Y$, let $\di^v\BY=\eta_i\cup \xi_{i+1}$ 
stand for the part of the $J$-boundary of $\BY$ attached to $v$.
%We let $\hat S_Y^v$ be the part of the $J$-boundary
%attached to the vertex $v$ of $Y$.

\comm{*********************
Note that the immersion 
constructed in Lemma \ref{psi-maps} extends continuously to the boundary of the 
puzzle piece $\BY$.
%  minus finitely many points where the $J$-boundary meets the $O$-boundery. 
(However, $\i(\di \BY)$ is not contained in $\di\BZ$, unless $Z=Y$.)  
In what follows we will assume this extension without further comment.


% In what follows it will be convenient to consider possibly disconnected curves.
A {\it multicurve} in some space $X$ is a continuous map $\gamma: \cup_{k=1}^l [s_k, t_k]\ra X$
parametrized by a finite union of disjoint intervals $[s_k, t_k]\subset \R$.% 
\footnote{We allow that the boundary points of a  multicurve in a $\psi$-puzzle piece $\BY$ belong to  $\di_J\BY$.} 
Note that  multicurves are ordered.
A multicurve in a puzzle piece $\BY$  is called {\it horizontal} if 
$$ 
   \gamma(s_1)\in \di^{v_0} \BY, \quad  \gamma(t_k), \gamma(s_{k+1})\in \di^{v_k}\BY,\ k=1, \dots, l-1, 
                  \quad \gamma(t_l)\in \di^{v_l} \BY
$$
for some vertices $v_k$ of $\BY$, $k=0, \dots, l$. We say that such a multicurve ``connects''
$\di^{v_0} \BY $ to $\di^{v_l} \BY$.
The following statement motivates introduction of  multicurves: 

\begin{lem}\label{lifts}
Let $v$ and $w$ be two vertices of a geometric puzzle piece $Y\in\YY^m(l)$. 
Then any curve $\gamma$ in $\BU^l$ connecting $S_Y^v$ to $S_Y^w$ 
contains a multicurve $\gamma'$  that lifts to a multicurve  $\gamma^*$ in $\BY$
connecting $\di_v\BY$ to $\di_w\BY$.
\end{lem}

 

Given two vertices $v$ and $w$ of $Y$, let $\GG_Y(v,w)$ stand for the family of horizontal multicurves in $\BY$
connecting $\di^v(\BY)$ to $\di^w(\BY)$. 
Finally, let 
\begin{equation}\label{Bd}
         \Bd_Y(v,w)= \LL(\GG_Y(v,w))
\end{equation}
stand for the extremal distance between the corresponding parts of $J$-boundary of $\BY$.
% Then we immediately conclude:

\begin{lem}\label{Bd}
  If  $f|Y$ is univalent,
then $\Bd_Y(v,w)=\Bd_{f(Y)}(fv, fw). $  
\end{lem}
************************************************}

\comm{****************

\subsection{Pseudo-summary}\label{psi-summary}

   We can now recast the Summary of \S \ref{summary} in the pseudo-context.  
We start with a $u\psi$-PL map $\bf$, 
restrict it to a unicritical PL map $f$ and construct the associated dynasty of kingdoms. 
We fix an arbitrary $m$, assume that the height of the dynasty is greater than $\log_2 m +5$,
and take some $n>\log_2 m +5$.
Then for any domain $\La =\La_k$, we obtain  the map $\Psi=\Psi_k$ \ref{3 domains}
satisfying properties (P1)-(P4). Applying to it the $\psi$-functor, we obtain:

\begin{itemize}

\item [(P1$'$)] An even-valued immersion ${\Bpsi}^{-1}: \BV^0 \ra \BUps$
 of degree bounded by $d^{2n+m}$;

\item[(P2$'$)] An even-valued immersion  ${\Bpsi}^{-1}: \BLa'\ra \BF$
  of degree bounded by  $d^5$;

\item [(P3$'$)] An immersion $\BUps\ra \BV^n$;   

\item[(P4$'$)] An immersion $\hat\BE^{i+1} \ra  \BLa' $  which restricts to 
   a homeomorphism $K_{E^{i+1}}\ra   K_\La$.

\end{itemize}

Moreover, 

\begin{itemize}
\item[(P5$'$)] The images of the  immersions
 $\BLa_k'\ra \BV$ are disjoint from the Julia pieces $K_{\La_j}$,  $k\not=j$.

\item[(P6$'$)]
There are at least $m/2$ domains $\La_k$.

\end{itemize}
******************************************}











